\documentclass{article}

\usepackage[utf8]{inputenc}
\usepackage[T1]{fontenc}
\usepackage{helvet}
\usepackage{geometry}
\usepackage{xcolor}
\usepackage{eso-pic}
\usepackage{fancyhdr}
\usepackage{parskip}
\usepackage{microtype}
\usepackage[english, ukrainian]{babel}
\usepackage{url}
\usepackage{amsmath}
\usepackage{amssymb}
\usepackage{amsthm}
\usepackage{longtable}
\usepackage{booktabs}
\usepackage{hyperref}
\usepackage{titlesec}
\usepackage{tocloft}

\renewcommand{\cftsecfont}{\small\bfseries}
\renewcommand{\cftsubsecfont}{\footnotesize}
\renewcommand{\cftsubsubsecfont}{\scriptsize}
\renewcommand{\cftsecpagefont}{\small\bfseries}
\renewcommand{\cftsubsecpagefont}{\footnotesize}
\renewcommand{\cftsubsubsecpagefont}{\scriptsize}

\definecolor{nato}{HTML}{293757}
\definecolor{footergray}{gray}{0.65}

\geometry{a4paper,left=3cm,right=3cm,top=3cm,bottom=3cm}
\renewcommand{\familydefault}{\sfdefault}
\setcounter{secnumdepth}{3}

\titleformat{\section}
  {\color{nato}\Huge\bfseries\raggedright}{\thesection.}{1em}{}
\titlespacing*{\section}{0pt}{30pt}{12pt}

\titleformat{\subsection}
  {\color{nato}\LARGE\bfseries\raggedright}{\thesubsection.}{1em}{}
\titlespacing*{\subsection}{0pt}{20pt}{8pt}

\titleformat{\subsubsection}
  {\color{nato}\Large\bfseries\raggedright}{\thesubsubsection.}{1em}{}
\titlespacing*{\subsubsection}{0pt}{10pt}{6pt}

\fancyhf{}
\fancyhead[R]{\small\color{footergray} 2 червня 2026}
\fancyfoot[L]{\small\color{footergray} Комплексна система захисту інформації. Цільовий профіль безпеки вищих судів (3).}
\fancyfoot[R]{\small\color{footergray}\thepage}

\newcommand\HUGE[1]{\fontsize{40}{40}\selectfont #1}

\renewcommand{\headrulewidth}{0pt}
\renewcommand{\footrulewidth}{0.4pt}
\renewcommand{\footrule}{\color{footergray}\hrule width \textwidth height 0.4pt}

\begin{document}

\pagestyle{fancy}

\begin{titlepage}
  \AddToShipoutPictureBG*{\AtPageLowerLeft{\color{nato}\rule{\paperwidth}{\paperheight}}}
  \color{white}\thispagestyle{empty}\vspace*{4cm}\centering
  {\Huge  \textbf{Комплексна система захисту інформації} \par} \vspace{2cm}
  {\HUGE  \textbf{Цільовий профіль безпеки вищих судів (3)} \par} \vspace{2.5cm}
  \vfill
  {\large 2026 \copyright\ Критпографічні Телесистеми \par}
\end{titlepage}

\newpage

\begin{center}
  {\Huge\color{nato}\bfseries Зміст\par}
\end{center}
\vspace{40pt}
\tableofcontents

\begin{abstract}
Цільовий профіль безпеки (ЦПБ) — індивідуальний для конкретної системи підприємства. Саме на його основі створюється/модернізується КЗЗІ. Включає лише додаткові вимоги (відмінності) відносно Галузевого профілю.
\end{abstract}

\section{CP}
Клас заходів захисту CP --- ПЛАНУВАННЯ БЕЗПЕРЕРВНОЇ РОБОТИ

Цей клас описує заходи для відновлення функціонування інформаційної системи та забезпечення її безперебійної роботи у разі надзвичайних ситуацій.

\textbf{Перелік заходів захисту:} Id-spe-cp-2; Id-spe-cp-10.

\subsection{id-spe-cp-2}


\vspace{10pt}
{\footnotesize\bfseries
No: 1\\
Name: cp\_\allowbreak{}2\_\allowbreak{}a\_\allowbreak{}01\\
Type: list\\
Default: [\textquotedbl{}admin\textquotedbl{}, \textquotedbl{}security\_\allowbreak{}officer\textquotedbl{}]
}

{\footnotesize Розроблено план забезпечення безперервної роботи у надзвичайних ситуаціях для системи, який визначає основні завдання, функції та пов'язані з ними вимоги щодо безперервної роботи; CP-02(a)[02][01] розроблено план забезпечення безперервної роботи у надзвичайних ситуаціях для системи, який забезпечує цілі; CP-02(a)[02][02] розроблено план забезпечення безперервної роботи у надзвичайних ситуаціях для системи, який забезпечує пріорітети; CP-02(a)[02][03] розроблено план забезпечення безперервної роботи у надзвичайних ситуаціях для системи, який забезпечує відповідні показники; CP-02(a)[03][01] розроблено план забезпечення безперервної роботи на випадок надзвичайних ситуацій для системи, в якому визначено ролі; CP-02(a)[03][02] розроблено план забезпечення безперервної роботи на випадок надзвичайних ситуацій для системи, в якому визначено обов'язки; CP-02(a)[03][03] розроблено план забезпечення безперервної роботи на випадок надзвичайних ситуацій для системи, в якому визначено відповідальних осіб з контактною інформацією}


{\footnotesize\bfseries
No: 2\\
Name: cp\_\allowbreak{}2\_\allowbreak{}a\_\allowbreak{}04\\
Type: string\\
Default: nil
}

{\footnotesize Розроблено план забезпечення безперервної роботи на випадок непередбачених обставин для системи, який спрямований на підтримку основних завдань і функції, попри системні збої, компрометації або помилки}


{\footnotesize\bfseries
No: 3\\
Name: cp\_\allowbreak{}2\_\allowbreak{}a\_\allowbreak{}05\\
Type: list\\
Default: [\textquotedbl{}admin\textquotedbl{}, \textquotedbl{}security\_\allowbreak{}officer\textquotedbl{}]
}

{\footnotesize Розроблено план забезпечення безперервної роботи у надзвичайних ситуаціях для системи, який спрямований на повне відновлення функціонування системи без погіршення запланованих і реалізованих заходів захисту інформації та персональних даних}


{\footnotesize\bfseries
No: 4\\
Name: cp\_\allowbreak{}2\_\allowbreak{}a\_\allowbreak{}06\\
Type: list\\
Default: [\textquotedbl{}admin\textquotedbl{}, \textquotedbl{}security\_\allowbreak{}officer\textquotedbl{}]
}

{\footnotesize Розроблено план забезпечення безперервної роботи на випадок надзвичайних ситуацій для системи, який вирішує питання обміну інформацією про надзвичайні ситуації; CP-02(a)[07][01] розроблено план забезпечення безперервної роботи у надзвичайних ситуаціях для системи, який переглядається персоналом або ролями; CP-02(a)[07][02] розроблено план забезпечення безперервної роботи у надзвичайних ситуаціях для системи, який затверджено персоналом або ролями}


{\footnotesize\bfseries
No: 5\\
Name: cp\_\allowbreak{}2\_\allowbreak{}b\_\allowbreak{}01\\
Type: list\\
Default: [\textquotedbl{}admin\textquotedbl{}, \textquotedbl{}security\_\allowbreak{}officer\textquotedbl{}]
}

{\footnotesize Копії плану забезпечення безперервної роботи на випадок надзвичайних ситуацій розповсюджуються серед персоналу}


{\footnotesize\bfseries
No: 6\\
Name: cp\_\allowbreak{}2\_\allowbreak{}b\_\allowbreak{}02\\
Type: string\\
Default: nil
}

{\footnotesize Копії плану забезпечення безперервної роботи на випадок надзвичайних ситуацій розповсюджуються серед елементів}


{\footnotesize\bfseries
No: 7\\
Name: cp\_\allowbreak{}2\_\allowbreak{}c\\
Type: string\\
Default: \textquotedbl{}щорічно\textquotedbl{}
}

{\footnotesize Діяльність з планування безперервної роботи координується з діяльністю із заходами по усуненню інцидентів; CP-02(d) переглядається план забезпечення безперервної роботи у надзвичайних ситуаціях для системи частота}


{\footnotesize\bfseries
No: 8\\
Name: cp\_\allowbreak{}2\_\allowbreak{}e\_\allowbreak{}01\\
Type: string\\
Default: nil
}

{\footnotesize План забезпечення безперервної роботи на випадок надзвичайних ситуацій оновлюється з урахуванням змін в організації, системі або середовищі функціонування}


{\footnotesize\bfseries
No: 9\\
Name: cp\_\allowbreak{}2\_\allowbreak{}e\_\allowbreak{}02\\
Type: integer\\
Default: 30
}

{\footnotesize План забезпечення безперервної роботи на випадок надзвичайних ситуацій оновлюється для вирішення проблем, що виникають під час впровадження, виконання або тестування плану дій на випадок надзвичайних ситуацій}


{\footnotesize\bfseries
No: 10\\
Name: cp\_\allowbreak{}2\_\allowbreak{}f\_\allowbreak{}01\\
Type: list\\
Default: [\textquotedbl{}admin\textquotedbl{}, \textquotedbl{}security\_\allowbreak{}officer\textquotedbl{}]
}

{\footnotesize Зміни в плані забезпечення безперервної роботи на випадок надзвичайних ситуацій повідомляються персоналу}


{\footnotesize\bfseries
No: 11\\
Name: cp\_\allowbreak{}2\_\allowbreak{}f\_\allowbreak{}02\\
Type: string\\
Default: nil
}

{\footnotesize Зміни в плані забезпечення безперервної роботи на випадок надзвичайних ситуацій повідомляються елементам}


{\footnotesize\bfseries
No: 12\\
Name: cp\_\allowbreak{}2\_\allowbreak{}g\_\allowbreak{}01\\
Type: integer\\
Default: 30
}

{\footnotesize Уроки, отримані під час тестування планів забезпечення безперервної роботи у надзвичайних ситуаціях або фактичних дій у надзвичайних ситуаціях, включаються в навчання}


{\footnotesize\bfseries
No: 13\\
Name: cp\_\allowbreak{}2\_\allowbreak{}g\_\allowbreak{}02\\
Type: integer\\
Default: 30
}

{\footnotesize Уроки, отримані під час тестування планів забезпечення безперервної роботи у надзвичайних ситуаціях або фактичних дій у надзвичайних ситуаціях, включаються в тестування}


{\footnotesize\bfseries
No: 14\\
Name: cp\_\allowbreak{}2\_\allowbreak{}h\_\allowbreak{}01\\
Type: string\\
Default: nil
}

{\footnotesize План забезпечення безперервної роботи у надзвичайних ситуаціях захищений від несанкціонованого доступу}


{\footnotesize\bfseries
No: 15\\
Name: cp\_\allowbreak{}2\_\allowbreak{}h\_\allowbreak{}02\\
Type: string\\
Default: nil
}

{\footnotesize План забезпечення безперервної роботи у надзвичайних ситуаціях захищений від несанкціонованих змін}


{\footnotesize\bfseries
No: 16\\
Name: cp\_\allowbreak{}2\_\allowbreak{}odp\_\allowbreak{}01\\
Type: list\\
Default: [\textquotedbl{}admin\textquotedbl{}, \textquotedbl{}security\_\allowbreak{}officer\textquotedbl{}]
}

{\footnotesize Визначено персонал або ролі для перегляду плану забезпечення безперервної роботи у надзвичайних ситуаціях}


{\footnotesize\bfseries
No: 17\\
Name: cp\_\allowbreak{}2\_\allowbreak{}odp\_\allowbreak{}02\\
Type: list\\
Default: [\textquotedbl{}admin\textquotedbl{}, \textquotedbl{}security\_\allowbreak{}officer\textquotedbl{}]
}

{\footnotesize Визначено персонал або ролі для затвердження плану забезпечення безперервної роботи у надзвичайних ситуаціях}


{\footnotesize\bfseries
No: 18\\
Name: cp\_\allowbreak{}2\_\allowbreak{}odp\_\allowbreak{}03\\
Type: list\\
Default: [\textquotedbl{}admin\textquotedbl{}, \textquotedbl{}security\_\allowbreak{}officer\textquotedbl{}]
}

{\footnotesize Визначено ключовий резервний персонал (ідентифікований за іменами та/або за ролями), якому поширюються копії плану забезпечення безперервної роботи на випадок надзвичайних ситуацій}


{\footnotesize\bfseries
No: 19\\
Name: cp\_\allowbreak{}2\_\allowbreak{}odp\_\allowbreak{}04\\
Type: string\\
Default: nil
}

{\footnotesize Визначено ключові елементи, на які поширюються копії плану забезпечення безперервної роботи на випадок надзвичайних ситуацій}


{\footnotesize\bfseries
No: 20\\
Name: cp\_\allowbreak{}2\_\allowbreak{}odp\_\allowbreak{}05\\
Type: string\\
Default: \textquotedbl{}щорічно\textquotedbl{}
}

{\footnotesize Визначено періодичність перегляду плану забезпечення безперервної роботи у надзвичайних ситуаціях}


{\footnotesize\bfseries
No: 21\\
Name: cp\_\allowbreak{}2\_\allowbreak{}odp\_\allowbreak{}06\\
Type: list\\
Default: [\textquotedbl{}admin\textquotedbl{}, \textquotedbl{}security\_\allowbreak{}officer\textquotedbl{}]
}

{\footnotesize Визначено ключовий резервний персонал (ідентифікований за іменами та/або ролями), якому необхідно повідомити про зміни}


{\footnotesize\bfseries
No: 22\\
Name: cp\_\allowbreak{}2\_\allowbreak{}odp\_\allowbreak{}07\\
Type: string\\
Default: nil
}

{\footnotesize Визначено ключові елементи організації, і які необхідно повідомити про зміни}




\subsection{id-spe-cp-10}


\vspace{10pt}
{\footnotesize\bfseries
No: 1\\
Name: cp\_\allowbreak{}10\_\allowbreak{}01\\
Type: integer\\
Default: 30
}

{\footnotesize Відновлення системи до відомого стану забезпечується протягом часу після збою, компрометації або помилки}


{\footnotesize\bfseries
No: 2\\
Name: cp\_\allowbreak{}10\_\allowbreak{}02\\
Type: integer\\
Default: 30
}

{\footnotesize Відтворення системи до відомого стану забезпечується протягом часу після збою, компрометації або помилки}


{\footnotesize\bfseries
No: 3\\
Name: cp\_\allowbreak{}10\_\allowbreak{}odp\_\allowbreak{}02\\
Type: integer\\
Default: 30
}

{\footnotesize Визначено період часу для відновлення, часу та цілям відновлення системи; визначено період часу для відтворення, часу та цілям відновлення системи}




\section{SC}
Клас заходів захисту SC --- ЗАХИСТ ІНФОРМАЦІЙНОЇ СИСТЕМИ ТА

Цей клас охоплює механізми захисту інформації під час її передачі мережами зв'язку, криптографічний захист та ізоляцію критичних компонентів.

\textbf{Перелік заходів захисту:} Доступність ресурсів (SC-6).

\subsection{ДОСТУПНІСТЬ РЕСУРСІВ (SC-6)}
Забезпечити захист доступності ресурсів, виділивши [Призначення: визначені організацією ресурси], по [Вибір (один або кілька); пріоритет; квоти; [Призначення: визначені організацією заходи з безпеки]].

\vspace{10pt}
{\footnotesize\bfseries
No: 1\\
Name: sc\_\allowbreak{}6\_\allowbreak{}01\\
Type: string\\
Default: nil
}

{\footnotesize Захищено доступність ресурсів шляхом розподілу ресурсів за ВИБІРКОВИМ ЗНАЧЕННЯМ ПАРАМЕТРА(ів)}


{\footnotesize\bfseries
No: 2\\
Name: sc\_\allowbreak{}6\_\allowbreak{}odp\_\allowbreak{}01\\
Type: string\\
Default: nil
}

{\footnotesize Визначені ресурси, які необхідно виділити для захисту доступності ресурсів}


{\footnotesize\bfseries
No: 3\\
Name: sc\_\allowbreak{}6\_\allowbreak{}odp\_\allowbreak{}02\\
Type: string\\
Default: nil
}

{\footnotesize Вибрано одне або декілька з наступних ЗНАЧЕНЬ ПАРАМЕТРІВ: \{priority; quota; controls\}}


{\footnotesize\bfseries
No: 4\\
Name: sc\_\allowbreak{}6\_\allowbreak{}odp\_\allowbreak{}03\\
Type: string\\
Default: \textquotedbl{}автоматизований засіб моніторингу\textquotedbl{}
}

{\footnotesize Визначені засоби контролю для захисту доступності ресурсів (якщо вибрано)}




\end{document}
